% Independent proof. Preserve the previously frozen edition unchanged.
\section{Quadratic reciprocity excludes both remaining positive Type I events}
This section settles the two residue events left open in PES9--10.
The proof uses the original Type I parameter identities of Christian
Elsholtz and Terence Tao, \emph{Counting the number of solutions to the
Erd\H{o}s--Straus equation on unit fractions}, arXiv1107.1010v6,
Section~2 \cite[Section~2]{ElsholtzTao2015},
Proposition \texttt{type-1} and equations \texttt{I-1}--\texttt{I-9},
and their displayed quadratic-reciprocity and supplementary laws
\texttt{quadratic}, \texttt{quadratic-2}. Their parameter order
$(a,b,c,d,e,f)$ is exactly our $(a,b,d,h,k,r)$, and their denominator
order $(x,y,z)$ is our $(z,x,y)$. Both substitutions are explicit:
no original ES point, sign, or coefficient is changed. The preserved
indexed author source has SHA256
\texttt{b0469a67a737b7f4e77778310c87e22f5783ae9704f70aa55919a4a40104611c}.
Its displayed \texttt{quadratic-1} prints the exponent $(n-1)/4$;
the correct law used here is $(-1/n)=(-1)^{(n-1)/2}$. This correction
is stated rather than silently attributed to the printed source.

\subsection{All parameters and all coprimalities}
Let $p\equiv1\pmod{12}$ be a rational prime and let a positive
one-divisible ES witness be marked, as in ISR, by
$(\mathrm M,\mathrm N,\mathrm T,\mathrm E)=(p,x,y,z)$ with $z=pZ$
and $p\nmid xyZ$. Choose the recorded permutation of the two unit
denominator labels for which $x<y$. Write $x=g a$, $y=g b$ with
$g=\gcd(x,y)$, $a<b$, $\gcd(a,b)=1$. Positivity of the reciprocal
equation gives $4gab-p(a+b)>0$. Put $d=4gab-p(a+b)$. Then
$Z=gab/d$, and $\gcd(d,a)=\gcd(d,b)=1$ since reducing $d$ modulo
$a$ and $b$ gives $-pb$ and $-pa$. Thus $d\mid g$; define $h=g/d$.
Also $p\nmid d$ since $d\equiv4gab\not\equiv0\pmod p$.
Rearranging gives $d(4hab-1)=p(a+b)$, so $p\nmid d$ proves
$p\mid4hab-1$. The following positive integers are consequently
defined without a divisibility assumption:
\begin{equation}
 \begin{aligned}
 &x=hda,\quad y=hdb,\quad z=phab,\quad
 k=(4hab-1)/p=(a+b)/d,\\
 &r=4x-p=4had-p>0,\qquad
 pk=4hab-1,\quad dk=a+b,\quad rk=4ha^2+1,\quad b=dk-a.
 \end{aligned}                                          \tag{OE1}
\end{equation}
For the third identity, substitute $b=dk-a$ into $pk=4hab-1$
and rearrange. This proves the same complete parameter identities
as the stated Elsholtz--Tao crosswalk. In particular
\begin{equation}
 p\nmid habd,\quad \gcd(h,k)=\gcd(a,k)=\gcd(b,k)=1,
 \quad k\equiv r\equiv3\pmod4,
 \quad pk\equiv-1\pmod{4h}.                            \tag{OE2}
\end{equation}
The first assertion follows from $p\nmid xyZ$.
Any common divisor of $h$ and $k$ would divide $4hab-pk=1$.
A common divisor of $a$ and $k$ would divide $a$ and $dk-a=b$,
and is one; the argument for $b$ is identical. Finally
$pk\equiv-1\pmod4$ and $p\equiv1\pmod4$ give $k\equiv3$;
$r=4had-p$ gives $r\equiv3$. These facts do not assume either
$p\nmid(x-y)$ or $p\nmid S$ and therefore do not assume the result
that will be proved.

\subsection{The exact Jacobi calculation}
For a positive odd integer $n$ and $\gcd(t,n)=1$, use the Jacobi
symbol defined as the product of the Legendre symbols at all prime
factors of $n$, counted with their multiplicities; the symbol with
denominator one is one. The exact classical laws used below are
\begin{equation}
 \left(\frac mn\right)\left(\frac nm\right)
    =(-1)^{(m-1)(n-1)/4},\qquad
 \left(\frac{-1}{n}\right)=(-1)^{(n-1)/2},\qquad
 \left(\frac2n\right)=(-1)^{(n^2-1)/8},                 \tag{OE3}
\end{equation}
for positive odd coprime $m,n$. These are the quadratic-reciprocity
law and its supplementary laws in the exact conventions just stated.
The corrected middle exponent can also be checked directly: for an
odd prime $q$, Euler's criterion gives $(-1/q)=(-1)^{(q-1)/2}$;
prime factorization and
$(uv-1)/2\equiv(u-1)/2+(v-1)/2\pmod2$ for odd $u,v$
give the same formula for every odd $n$. Thus the typographical
exponent in the indexed author file is never used as a premise.

Reduction of $rk=4ha^2+1$ modulo $k$ gives $4ha^2=-1\pmod k$.
The coprimalities in OE2 imply that every Jacobi symbol below is
$+1$ or $-1$, not zero. Since four and $a^2$ are invertible squares,
\begin{equation}
 \left(\frac hk\right)
    =\left(\frac{-1}{k}\right)=-1.                    \tag{OE4}
\end{equation}
This remains valid when $k$ is composite; no primality of $k$ has
been assumed.

Write $h=2^e h_0$ with $h_0$ positive odd. Reciprocity, with
$p\equiv1\pmod4$ and $k\equiv3\pmod4$, gives
\begin{align}
 \left(\frac{h_0}{p}\right)\left(\frac{h_0}{k}\right)
 &=\left(\frac p{h_0}\right)\left(\frac k{h_0}\right)
       (-1)^{(h_0-1)/2} \nonumber\\
 &=\left(\frac{pk}{h_0}\right)(-1)^{(h_0-1)/2}
  =\left(\frac{-1}{h_0}\right)(-1)^{(h_0-1)/2}=1.
                                                               \tag{OE5}
\end{align}
Here $pk\equiv-1\pmod{h_0}$ is the unchanged OE2 relation.
When $h_0=1$, all its symbols and the sign are one, so this case
is included. If $e=0$ there is no power of two left to check.
If $e>0$, then $pk=4hab-1\equiv-1\pmod8$. The last law in OE3
and multiplicativity in its positive odd denominator give
\begin{equation}
 \left(\frac2p\right)\left(\frac2k\right)
     =\left(\frac2{pk}\right)=1.
 \qquad
 \left(\frac hp\right)\left(\frac hk\right)=1
 \quad\text{in both parity cases}.                     \tag{OE6}
\end{equation}
Combining OE4 and OE6 proves the exact nonresidue statement
\begin{equation}
 \boxed{\left(\frac hp\right)=-1.}\qquad
 \left(\frac{ab}{p}\right)=-1,
 \qquad \boxed{\left(\frac xp\right)
                       \left(\frac yp\right)=-1.}      \tag{OE7}
\end{equation}
For the middle equality, reduce $4hab=pk+1$ modulo $p$ and use
that four is a square. For the final equality use
$xy=h^2d^2ab$, with $hd$ a $p$-unit. Thus the two original unit
denominators occupy opposite quadratic-residue classes.

\subsection{Both previously remaining events are excluded}
If $x\equiv y\pmod p$, their Legendre symbols are equal, contrary
to OE7. If $x\equiv-y\pmod p$, their symbols are again equal
because $(-1/p)=1$ for $p\equiv1\pmod4$. Both events are therefore
impossible. Since $S=p+x+y+pZ\equiv x+y\pmod p$, this proves
\begin{equation}
 \boxed{v_p(x-y)=0,\qquad v_p(S)=0.}                    \tag{OE8}
\end{equation}
The result is stronger than the earlier positivity bound
$v_p(x-y)\le1$, and it uses no search cutoff. Sorting only selected
which unit denominator is called $x$ during the proof. OE7--8 are
symmetric under exchanging the two original unit labels, so they
return to the original marking without a change of conclusion.

In the two-divisible case $S\equiv x\not\equiv0\pmod p$ was
already immediate, and PES5--7 excluded every closer scaled pair.
Consequently \emph{every} actual positive witness in the present
hard-prime class has the original coefficient $A=-1/S$ a $p$-unit.
ISR18--23 remains a valid calculation of nonintegral charts for more
general quartics, but those charts do not occur in this actual
positive witness class. No coordinate shift or quadratic twist is
needed to make an actual witness integral.

\subsection{Exactly two actual eight-state integral strata}
Keep the four original roots $(p,x,y,z)$, $D_i=Af'(r_i)$,
$u_i=p^{-b_i}D_i$, and the full matrices $C_p,G_I$ from ISR6,12.
All units in those matrices are unchanged. Combining OE8 with
PES7 gives the exhaustive actual lists
\begin{equation}
 \begin{array}{c|c|c|c|c|c}
 & (b_i)&(m_i)&(\epsilon_i)&(v_pI_i)&
       \operatorname{Smith}(O/E)=\operatorname{Smith}(E/I)\\\hline
 \text{one}&(1,0,0,1)&(0,0,0,0)&(1,0,0,1)&(1,0,0,1)&
                  (0,0,0,0,0,0,1,1)\\
 \text{two}&(2,0,2,2)&(1,0,1,1)&(0,0,0,0)&(3,0,3,3)&
                  (0,0,0,1,1,2,2,3).
 \end{array}                                            \tag{OE9}
\end{equation}
In the first row $z$ is the divisible denominator; in the second
row $y,z$ are. Any original marking uses the corresponding recorded
permutation. The first row follows because the only positive gap
valuation is $v_p(p-z)=1$ from ISR2: the other four mixed gaps are
units by their two different residue classes, and $x-y$ is a unit
by OE8. The second row is exactly PES7. Thus
\begin{equation}
 \ell(O/E)=\ell(E/I)=\begin{cases}2&\text{one},\\9&\text{two},\end{cases}
 \qquad \ell(O/I)=\begin{cases}4&\text{one},\\18&\text{two}.
                      \end{cases}                      \tag{OE10}
\end{equation}
The complete eight generators remain $Ap^{m_i}f_i$ and
$A\sigma f_i$; their branch ideals are still
$I_i=p^{m_i}D_iO_i$, with the full $D_i$ rather than a selected
unit representative.

The original quartic coefficient is $C_0=-5Axyz$, with valuation
one or two in the respective rows. The exact specialization of
the rank-twenty-four coefficient algebra proved in PS14 now gives
\begin{equation}
 \Pi_P(C_0^2E_N)=C_0^2E\subset I
 \quad\text{for every actual positive witness},\qquad
 \min\{j\ge0:C_0^jE\subset I\}
  =\begin{cases}1&\text{one},\\2&\text{two}.
                      \end{cases}                      \tag{OE11}
\end{equation}
Unlike the earlier formulation, this conclusion has no unverified
$p\nmid S$ restriction for the stated actual witness class: OE8
and the two-divisible observation prove it in every case.

There is a full calculation of the intermediate quotient, not just
the inclusion. If the eight Smith exponents in OE9 are $e_i$ and
$c=v_p(C_0)$, then for every $j$ satisfying OE11,
\begin{equation}
 I/(C_0^jE)\simeq\bigoplus_{i=1}^8R/(p^{jc-e_i}),
 \qquad
 Z_j=G_I^{-1}C_0^jI_8
  =\operatorname{diag}\left(
   \operatorname{diag}\left(\frac{C_0^j}{p^{m_i}D_i}\right)V,
   \operatorname{diag}\left(\frac{C_0^j}{D_i}\right)V\right).
                                                               \tag{OE12}
\end{equation}
To prove this, choose the integral invertible left and right
matrices expressing the Smith form of $G_I$. Then $G_I^{-1}C_0^j$
has the complementary diagonal exponents $jc-e_i$, with the full
unit of $C_0^j$ multiplying an invertible matrix. All these
exponents are nonnegative by the already proved inclusion. The
explicit matrix is PES11 and retains every original entry. Thus
the actual quotient for the coefficient conductor has, in
increasing exponent order,
\begin{equation}
 \begin{array}{c|c|c}
 & \operatorname{Smith}\bigl(I/(C_0^2E)\bigr)&
                 \ell\bigl(I/(C_0^2E)\bigr)\\\hline
 \text{one}&(1,1,2,2,2,2,2,2)&14\\
 \text{two}&(1,2,2,3,3,4,4,4)&23.
 \end{array}
 \qquad I/(C_0E)\simeq(R/pR)^6\quad\text{in the first row}.
                                                               \tag{OE13}
\end{equation}
Zero exponents in the last assertion are omitted as zero modules.
This gives the exact surviving difference between the two algebra
conductor ideals after selecting the actual witness.

\subsection{All original residue units and the complete reduced fibre}
For the one-divisible row, reducing the literal derivatives gives
\begin{equation}
 \begin{aligned}
 D_p/p&\equiv-\frac{3xy}{4(x+y)},&
 D_z/p&\equiv\frac{3xy}{4(x+y)},\\
 D_x&\equiv-\frac{x^2(x-y)}{x+y},&
 D_y&\equiv\frac{y^2(x-y)}{x+y}\pmod p.
 \end{aligned}                                          \tag{OE14}
\end{equation}
These are obtained by multiplying the three original root
differences, using $Z\equiv1/4$ and $A\equiv-1/(x+y)$.
Every denominator is invertible by OE8. Since $-1$ is a square
modulo $p$, the two ramified branches have the same quadratic
field class and the two unit branches have the same class:
the latter are both split or both unramified according as
$((x-y)/(x+y)\mid p)$ is $+1$ or $-1$.
This does not replace any full unit by its residue. For example,
an isomorphism between the two ramified generic factors uses the
actual lifted unit $c_{z,p}$ with $c_{z,p}^2=D_z/D_p$ and a chosen
residue root of $-1$, sending $\sigma_z\mapsto c_{z,p}\sigma_p$.
The elementary odd-prime lifting recurrence in ISR5 constructs
this unit and its negative. Both factor signs remain distinct.
The exterior cofactor in the stated Elsholtz--Tao crosswalk is
$D_E=k=(a+b)/d$; ISR30--31 therefore identify each ramified field
with $\mathbb Q_p(\sqrt{pk})$ using its actual lifted scalar.

For the entire one-divisible reduced fibre let bars denote
reduction in $\mathbb F_p$. Then
\begin{equation}
 \begin{aligned}
 E/pE\simeq{}&
 \mathbb F_p[r,\sigma]/(r^2,\sigma^2-2\bar A\bar x\bar y\,r)
 \\&\ \times\mathbb F_p[\sigma]/(\sigma^2-\bar D_x)
 \times\mathbb F_p[\sigma]/(\sigma^2-\bar D_y).
 \end{aligned}                                          \tag{OE15}
\end{equation}
The root factors $r^2,r-\bar x,r-\bar y$ are pairwise coprime
by OE8, so the displayed map is the full Chinese-remainder map
on both parity components. In the first factor put
$g_0=\bar A\bar x\bar y\ne0$. The original polynomial is
$\bar A r^2(r-\bar x)(r-\bar y)$ and its cofactor's linear
coefficient is $-\bar A(\bar x+\bar y)=1$. In the exact local
basis $(1,\sigma,\sigma^2,\sigma^3)$, the original residue and
trace have Gram matrices
\begin{equation}
 B_0=\begin{pmatrix}
 0&-g_0^{-2}&0&2\\-g_0^{-2}&0&2&0\\0&2&0&0\\2&0&0&0
 \end{pmatrix},\quad \det B_0=16,\qquad
 T_0=\operatorname{diag}(4,0,0,0).                     \tag{OE16}
\end{equation}
Indeed $\sigma^2=2g_0r$, and the original local residue gives
$\Lambda(\sigma)=-g_0^{-2}$, $\Lambda(\sigma r)=g_0^{-1}$,
$\Lambda(\sigma^3)=2$. All higher powers vanish, proving every
entry. The two other quadratic factors are separable because
their full reduced derivatives are nonzero and $p\ne2$; their
trace forms each have rank two. Thus the entire reduced trace
rank in this row is five, while the full eight-state residue
pairing remains perfect. In the two-divisible row the length-six
cluster has trace rank one and the other separable quadratic
factor rank two, so the entire reduced trace rank is three.
The original residue determinant $A^{-8}$ is a unit in both rows.

\subsection{The unchanged original coordinates and weighted conductor return}
For every row and root $r_i$, take both $\sigma_{i,\pm}$ satisfying
$\sigma_{i,\pm}^2=D_i$, and retain every original inverse coordinate
\begin{equation}
 q_{i,\pm}=\left(
 \sigma_{i,\pm}^{-1},\ -r_i-i\sigma_{i,\pm},\
 A\sigma_{i,\pm}^3+2r_i\sigma_{i,\pm}+3i\sigma_{i,\pm}^2,\
 7ir_i^2\sigma_{i,\pm}
 +(B-17r_i+Ar_i^2)\sigma_{i,\pm}^2
 -13i\sigma_{i,\pm}^3-2A\sigma_{i,\pm}^4\right).
                                                               \tag{OE17}
\end{equation}
These are precisely FS13, independently checked in ISR27; their
images under the original polynomial are the unchanged
$(A,B,C_0,D_0)$. OE7--16 change no coordinate, root label,
factor sign, original period, centre, or Gamma mass.

For the original fixed receiving maps $F_U,F_W$ of GD14 and ISR32,
the full new ideal factorization and its actual Gram matrices are
\begin{equation}
 \begin{aligned}
 &(F_U G_I)Z_j=C_0^jF_U,\qquad (F_W G_I)Z_j=C_0^jF_W,\\
 &(T_{\mathcal A}\oplus T_{\mathcal A})F_U G_I=F_WG_I,
 \qquad \Gamma_b^I=G_I^*\Gamma_b G_I\quad(b=U,W).
 \end{aligned}                                          \tag{OE18}
\end{equation}
Each equality follows by multiplying OE12 or the original receiver
identity. The determinant remains
$\det\Gamma_b^I=|A|^{16}|\det V|^4p^{2\sum m_i}\det\Gamma_b$,
with every original moment and mass in $\Gamma_b$ retained.
Thus the exclusion of the two remaining residue events makes the
arithmetic specialization and its original-metric return exhaustive
for actual positive witnesses. It does not identify the algebra
conductor ideal with the original linear conductor or assert an
RH endpoint from the finite arithmetic calculation.
